% GENERATED by tools/gen.py from data/vocab/u10.tsv -- do not edit.
\begin{vocabtab}
\vrow{\artf{die} Wohnung, Wohnungen}{mieszkanie}{CH: też die Wohnung}
\vrow{\artn{das} Haus, Häuser}{dom}{}
\vrow{\artn{das} Zimmer, Zimmer}{pokój}{}
\vrow{\artf{die} Küche, Küchen}{kuchnia}{}
\vrow{\artn{das} Bad, Bäder}{łazienka}{też: das Badezimmer}
\vrow{\artn{das} Schlafzimmer, Schlafzimmer}{sypialnia}{}
\vrow{\artn{das} Wohnzimmer, Wohnzimmer}{salon}{}
\vrow{\artm{der} Balkon, Balkone}{balkon}{}
\vrow{\artm{der} Flur, Flure}{korytarz}{}
\vrow{\artm{der} Keller, Keller}{piwnica}{}
\vrow{\artm{der} Garten, Gärten}{ogród}{}
\vrow{\artn{das} Bett, Betten}{łóżko}{}
\vrow{\artn{das} Sofa, Sofas}{sofa}{}
\vrow{\artm{der} Schrank, Schränke}{szafa}{}
\vrow{\artn{das} Regal, Regale}{regał}{}
\vrow{\artf{die} Miete, Mieten}{czynsz}{CH: der Mietzins}
\vrow{\artf{die} Nebenkosten}{opłaty dodatkowe}{tylko liczba mnoga}
\vrow{\artm{der} Vermieter, Vermieter}{wynajmujący}{f. die Vermieterin}
\vrow{\artf{die} Besichtigung, Besichtigungen}{oglądanie (mieszkania)}{}
\vrow{\artf{die} Kaution, Kautionen}{kaucja}{CH: die Mietkaution}
\vrow{\artn{das} Zentrum, Zentren}{centrum}{}
\vrow{\artm{der} Bahnhof, Bahnhöfe}{dworzec}{}
\vrow{\artm{der} Park, Parks}{park}{}
\vrow{\artf{die} Post}{poczta}{bez liczby mnogiej}
\vrow{\artf{die} Apotheke, Apotheken}{apteka}{}
\vrow{\artn{das} Krankenhaus, Krankenhäuser}{szpital}{CH: das Spital}
\vrow{\artf{die} Bank, Banken}{bank}{ława = die Bank, Bänke}
\vrow{in}{w, do}{Wo? in + Dativ; Wohin? in + Akkusativ}
\vrow{auf}{na}{}
\vrow{an}{przy, na (pionowo)}{}
\vrow{unter}{pod}{}
\vrow{über}{nad, ponad}{}
\vrow{neben}{obok}{}
\vrow{vor}{przed}{}
\vrow{hinter}{za}{}
\vrow{zwischen}{między}{}
\vrow{hier}{tutaj}{}
\vrow{da}{tam, tu}{}
\vrow{dort}{tam}{}
\vrow{links}{po lewej, w lewo}{}
\vrow{rechts}{po prawej, w prawo}{}
\vrow{oben}{na górze}{}
\vrow{unten}{na dole}{}
\vrow{zu Hause}{w domu}{}
\vrow{nach Hause}{do domu}{}
\vrow{stehen}{stać}{}
\vrow{liegen}{leżeć}{}
\vrow{sitzen}{siedzieć}{}
\vrow{mieten}{wynajmować (jako najemca)}{}
\vrow{gehen}{iść}{}
\vrow{bleiben}{zostawać}{}
\vrow{Wo ist das Büro?}{Gdzie jest biuro?}{}
\vrow{Wohin gehst du?}{Dokąd idziesz?}{}
\vrow{Ich wohne zur Miete.}{Mieszkam w wynajętym mieszkaniu.}{}
\vrow{Die Wohnung hat drei Zimmer.}{Mieszkanie ma trzy pokoje.}{}
\end{vocabtab}
